\section{导数的计算}

\subsection{导数的示例推导}

\begin{theorem}[二项式定理展开]
    对于任意正整数 $n$，二项式定理表示为：
    $$(a+b)^n = C_n^0 a^n + C_n^1 a^{n-1}b + C_n^2 a^{n-2}b^2 + \dots + C_n^k a^{n-k}b^k + \dots + C_n^n b^n$$
    同理，我们将 $a=x$, $b=h$ 代入二项式定理，得到 $(x+h)^n$ 的展开式：
    $$(x+h)^n = x^n + C_n^1 x^{n-1}h + C_n^2 x^{n-2}h^2 + \dots + C_n^n h^n$$
    $$(x+h)^n = x^n + nx^{n-1}h + \frac{n(n-1)}{2}x^{n-2}h^2 + \dots + h^n$$
\end{theorem}

\begin{solution}
    根据导数的定义 $f'(x) = \lim_{h \to 0} \frac{f(x + h) - f(x)}{h}$：
    \begin{align*}
        f'(x) & = \lim_{h \to 0} \frac{(x + h)^n - x^n}{h}                                                       \\
              & = \lim_{h \to 0} \frac{\left( x^n + nx^{n-1}h + C_n^2 x^{n-2}h^2 + \dots + h^n \right) - x^n}{h} \\
              & = \lim_{h \to 0} \frac{nx^{n-1}h + C_n^2 x^{n-2}h^2 + \dots + h^n}{h}
    \end{align*}
    因为 $h \to 0$ 但 $h \ne 0$，我们可以将分子中的 $h$ 因子提出来，并与分母的 $h$ 消去：
    \begin{align*}
         & = \lim_{h \to 0} \left( nx^{n-1} + C_n^2 x^{n-2}h + \dots + h^{n-1} \right)
    \end{align*}
    因此，当 $h \to 0$ 时，我们得到：$f'(x) = nx^{n-1}$
\end{solution}

\subsection{求导法则}

\begin{center}
    \begin{tabular}{l l l}
        \toprule
        \multicolumn{3}{c}{\textbf{基本初等函数}}                                                         \\
        \midrule
        $f(x) = c$            & $f'(x) = 0$                             & $c$ 为常数                   \\
        $f(x) = x^n$          & $f'(x) = nx^{n-1}$                      & $n \in \mathbb{R}$        \\
        $f(x) = \sqrt{x}$     & $f'(x) = \frac{1}{2\sqrt{x}}$           & $n = 1/2$ 的特例             \\
        $f(x) = \frac{1}{x}$  & $f'(x) = -\frac{1}{x^2}$                & $n = -1$ 的特例              \\
        \midrule
        \multicolumn{3}{c}{\textbf{指数与对数函数}}                                                        \\
        \midrule
        $f(x) = a^x$          & $f'(x) = a^x \ln a$                     & $a>0, a \neq 1$           \\
        $f(x) = \mathrm{e}^x$ & $f'(x) = \mathrm{e}^x$                  & $a=\mathrm{e}$ 的特例        \\
        $f(x) = \log_a x$     & $f'(x) = \frac{1}{x \ln a}$             & $a>0, a \neq 1, x>0$      \\
        $f(x) = \ln x$        & $f'(x) = \frac{1}{x}$                   & $a=\mathrm{e}$ 的特例, $x>0$ \\
        \midrule
        \multicolumn{3}{c}{\textbf{三角函数}}                                                           \\
        \midrule
        $f(x) = \sin x$       & $f'(x) = \cos x$                        &                           \\
        $f(x) = \cos x$       & $f'(x) = -\sin x$                       &                           \\
        $f(x) = \tan x$       & $f'(x) = \sec^2 x = \frac{1}{\cos^2 x}$ &                           \\
        \bottomrule
    \end{tabular}
\end{center}

除了基本函数的求导，还需要记住求导的运算法则：

\begin{center}
    \begin{tabular}{l l p{6cm}}
        \toprule
        \textbf{法则名称}                                                                & \textbf{公式} & \textbf{说明} \\
        \midrule

        \textbf{和差法则}                                                                &
        $[f(x) \pm g(x)]' = f'(x) \pm g'(x)$                                         &
        简单易懂                                                                                                     \\

        \addlinespace % 在行之间增加一点额外空间

        \textbf{乘法法则}                                                                &
        $[f(x)g(x)]' = f'(x)g(x) + f(x)g'(x)$                                        &
        前导后不导 + 前不导后导                                                                                            \\

        \addlinespace

        \textbf{商法则}                                                                 &
        $\left[\dfrac{f(x)}{g(x)}\right]' = \dfrac{f'(x)g(x) - f(x)g'(x)}{[g(x)]^2}$ &
        上导下不导 - 上不导下导                                                                                            \\

        \addlinespace

        \textbf{链式法则}                                                                &
        $[f(g(x))]' = f'(g(x)) \cdot g'(x)$                                          &
        复合函数，乘以再求一次。                                                                                             \\

        \bottomrule
    \end{tabular}
\end{center}

\v{2em}

\begin{example}
    下列函数求导正确的是（\quad ）
    \begin{itemize}
        \item 已知 $f(x)=x \ln x+x$ ，则 $f^{\prime}(x)=2+\ln x$
        \item 已知 $f(x)=\mathrm{e}^{2 x}$ ，则 $f^{\prime}(x)=2 \mathrm{e}^{2 x}$
        \item 已知 $f(x)=\frac{1}{x}$ ，则 $f^{\prime}(x)=\ln x$
        \item 已知 $f(x)=x^3+\sin x$ ，则 $f^{\prime}(x)=3 x^2+\cos x$
    \end{itemize}
\end{example}

\v{2em}

\begin{example}
    下列函数的导数计算正确的是（ \quad）
    \begin{itemize}
        \item $f(x)=x \mathrm{e}^x, f^{\prime}(x)=(x+1) \mathrm{e}^x$
        \item $f(x)=\sqrt{x}+1, f^{\prime}(x)=\frac{\sqrt{x}}{x}$
        \item $f(x)=\sin 2 x, f^{\prime}(x)=2 \cos 2 x$
        \item $f(x)=(2 x+1)^5, f^{\prime}(x)=10(2 x+1)^4$
    \end{itemize}
\end{example}

\subsection{含有导数值的函数求导}

\begin{theorem}
    只要将导数看成一个常数即可，后续解方程即可。
\end{theorem}

\begin{example}
    若函数 $f(x) = 6x^{3} - 2f'(1)x$，则 $f'(1) =$ （\quad）
    \begin{tasks}(4)
        \task 3 \task 6 \task $- 6$ \task $- 3$
    \end{tasks}
\end{example}

\v{8em}

\begin{example}
    已知函数 $f(x)=3 \ln (x-1)-x^2 f^{\prime}(2)$ ，则 $f^{\prime}(2)=$ \underline{$\qquad$} ．
\end{example}

\v{8em}

\begin{example}
    已知函数 $f(x)=\frac{1}{2} f^{\prime}(1) x^2+\ln x+\frac{f(1)}{3 x}$ ，则 $f^{\prime}(2)=$ \underline{$\qquad$} ．
\end{example}

\begin{problem}
如果要把上一题的系数 $\frac{1}{2} $ 改为 $1$ 呢？
\end{problem}

\v{8em}

\subsection{利用导数求切线方程}

\begin{theorem}
    \textbf{求切线方程}:
    \begin{itemize}
        \item \textbf{在点 $P(x_0, y_0)$ 处的切线}: P在曲线上，斜率为 $k=f'(x_0)$，使用点斜式 $y-y_0 = k(x-x_0)$。
        \item \textbf{过点 $P(x_0, y_0)$ 的切线}: P不在曲线上，需设切点为 $(x_1, f(x_1))$，表示出切线方程，再将点P坐标代入求解。
    \end{itemize}
\end{theorem}
\begin{example}
    已知函数 $f(x) = x^{2}$.
    \begin{enumerate}
        \item[(1)] 求曲线 $y = f(x)$ 在点 $(2,f(2))$ 处的切线方程；
        \item[(2)] 求曲线 $y = f(x)$ 过点 $(2,0)$ 的切线方程.
    \end{enumerate}
\end{example}

\v{8em}

\begin{example}
    已知曲线 $y = f(x)$ 在 $x = 5$ 处的切线方程是 $y = - 2x + 8$，则 $f(5) + f'(5)$ 的值为（\quad）
    \begin{tasks}(4)
        \task $- 4$ \task $- 2$ \task $2$ \task $4$
    \end{tasks}
\end{example}

\v{8em}

\begin{example}
    已知函数 $f(x) = 1 - 2x - \sin x$，则曲线 $y = f(x)$ 在点 $(0,f(0))$ 处的切线方程为（\quad）
    \begin{tasks}(4)
        \task $3x + y - 1 = 0$
        \task $2x - y + 1 = 0$
        \task $3x - y + 1 = 0$
        \task $2x + y - 1 = 0$
    \end{tasks}
\end{example}

\v{8em}

\begin{example}
    曲线 $f(x)=\frac{\ln x}{x^2}+\frac{1}{x+1}$ 在 $(1, f(1))$ 处的切线方程为
    \underline{$\qquad$} ．
\end{example}

\v{8em}

\begin{example}
    曲线 $y=\left(x^2-2\right) \ln x+x$ 在点 $(1,1)$ 处的切线方程为 \underline{$\qquad$} ．
\end{example}

\v{8em}

\begin{example}
    已知 $f(x)$ 为偶函数，曲线 $y = f(x)$ 在点 $(2,f(2))$ 处的切线与直线 $4x - y - 3 = 0$ 垂直，设 $f(x)$ 的导函数为 $f'(x)$，则 $f'(-2) = \underline{\hspace{3cm}}$.
\end{example}

\v{8em}

\begin{example}
    已知直线 $y=a x$ 是曲线 $y=\ln x$ 的切线，则实数 $a=$ \underline{$\qquad$}.
\end{example}

\v{8em}

\begin{example}
    已知函数 $f(x)=a x^2-2$ 的图象在点 $(\sqrt{3}, f(\sqrt{3}))$ 处的切线的倾斜角为 $\frac{\pi}{3}$ ，则曲线 $y=f(x)$ 在点 $(1, f(1))$ 处的切线的方程为（\quad）
    \begin{tasks}(4)
        \task $2 x+2 y-5=0$
        \task $2 x-2 y-5=0$
        \task $2 x-2 y+1=0$
        \task $2 x+2 y+1=0$
    \end{tasks}
\end{example}

\begin{example}
    过点 $(1,6)$ 且与曲线 $y=2 x^3+4$ 相切的直线方程为（\quad）
    \begin{itemize}
        \item $6 x-y=0$
        \item $9 x-y-3=0$
        \item $6 x-y=0$ 或 $3 x-2 y+9=0$
        \item $9 x-y-3=0$ 或 $3 x+2 y-15=0$
    \end{itemize}
\end{example}

\v{8em}

\begin{example}
    已知函数 $f(x) = x^{3} + x - 16$.
    \begin{enumerate}
        \item[(1)] 直线 $l$ 为曲线 $y = f(x)$ 在点 $(2, -6)$ 处的切线，求直线 $l$ 的方程；
        \item[(2)] 求过原点且与曲线 $y = f(x)$ 相切的切线方程及切点坐标．
    \end{enumerate}
\end{example}

\v{8em}
